El objetivo específico de esta tesina es desarrollar un sistema VI-SLAM que supere las limitaciones de los enfoques actuales en términos de precisión, robustez y eficiencia computacional, con énfasis en su aplicabilidad práctica sobre plataformas con recursos acotados. Para ello, se busca un sistema capaz de mantener una estimación de pose precisa frente a condiciones desafiantes, tales como escenas de baja textura, movimientos bruscos y trayectorias de larga duración, operando en tiempo real sin requerir hardware especializado. A fin de contener el error acumulado y garantizar consistencia a lo largo del tiempo, se prevé la incorporación de un mapa persistente con loop closure. Por último, el desempeño del sistema será evaluado exhaustivamente sobre datasets con amplia variedad de escenarios representativos de robótica y VR/AR, midiendo precisión y rendimiento frente a alternativas del estado del arte.
